\chapter{Исследовательский раздел}

В рамках лабораторной работы проведено сравнительное тестирование трёх больших языковых моделей (БЯМ) различных архитектур и производителей. Целью исследования являлось выявление способности моделей решать прикладные задачи разработки программного обеспечения (на языке Python с использованием библиотек `PyPDF2` и `re`) в условиях неидеально сформулированного технического задания.

\section{Объекты исследования}

В исследовании участвовали следующие модели:
\begin{enumerate}
    \item \textbf{Gemini 3 Pro Preview (Google)} --- облачная модель с поддержкой рассуждений (режим Thinking/Reasoning), позиционируемая как передовое решение (SOTA);
    \item \textbf{Qwen 3 Coder 30B Instruct (Alibaba)} --- открытая (Open Source) модель, развёрнутая локально и специализированная на написании программного кода;
    \item \textbf{YandexGPT 5.1 Pro (Яндекс)} --- отечественная облачная модель, используемая для оценки уровня импортозамещающих технологий в задачах кодинга.
\end{enumerate}

\section{Ход исследования и анализ взаимодействия}

Всем моделям подавался одинаковый входной запрос (промпт), содержащий ``сырую'' формулировку задания из методических указаний. Это позволило оценить не только навыки программирования, но и способность моделей интерпретировать технические требования.

\subsection{Gemini 3 Pro Preview}
\textbf{Первичная генерация.} Модель активировала блок размышлений, но выбрала неоптимальную стратегию:
\begin{itemize}
    \item \textit{алгоритм} --- использован метод \texttt{text.split(``$\backslash n$'')}, что делает невозможным поиск словосочетаний, разорванных переносом строки (например, ``рисунок'' на одной строке, а номер --- на другой);
    \item \textit{регулярные выражения} --- сразу корректно разделены ключевые слова ``от'' и ``до'' как независимые сущности.
\end{itemize}

\textbf{Итерация исправлений.} После указания на ошибку логики разбиения модель успешно переписала алгоритм. Финальное решение использует поиск по полному тексту и вычисление координат через подсчёт символов переноса, а также вводит список ``разрывающих'' символов (\texttt{[ $\backslash t\backslash n\backslash r\backslash f\backslash v$]+}).

\subsection{Qwen 3 Coder 30B (Local)}
\textbf{Первичная генерация.} Модель выдала синтаксически корректный код, но допустила ряд логических ошибок:
\begin{itemize}
    \item \textit{логика задачи} --- инвертировано условие (производился поиск наличия неразрывных пробелов вместо их отсутствия);
    \item \textit{алгоритм} --- использован ошибочный подход с \texttt{split(``$\backslash n$'')};
    \item \textit{регулярные выражения} --- предлоги ``от'' и ``до'' объединены в единую конструкцию, что приводит к игнорированию случаев, когда они встречаются по отдельности.
\end{itemize}

\textbf{Итерация исправлений.} Для получения рабочего результата потребовались уточнения. Модель с трудом восприняла указание разделить проверку предлогов ``от'' и ``до'' на независимые условия, потребовав явного переписывания запроса (промпта).

\subsection{YandexGPT 5.1 Pro}
\textbf{Первичная генерация.} Модель продемонстрировала противоречивый результат, предложив корректную алгоритмическую идею при неработоспособном программном коде:
\begin{itemize}
    \item \textit{алгоритм} --- это единственная модель, которая \textbf{сразу} предложила верный способ вычисления номера строки через \texttt{text[:match.start()].count(``$\backslash n$'') + 1}, работая со всем текстом целиком, а не с разбитым списком;
    \item \textit{качество кода} --- допущена грубая ошибка области видимости (Scope Error), при которой переменные, инициализированные внутри блока \texttt{with open(...) as file}, вызывались уже после закрытия файла, что приводило к аварийному завершению программы.
\end{itemize}

\textbf{Итерация исправлений.} Процесс отладки выявил низкую способность к самокоррекции. Модель дважды проигнорировала замечание об ошибке закрытого файла и, аналогично Qwen, испытывала трудности с разделением логики для предлогов ``от'' и ``до'' без кардинального изменения текста запроса.

\section{Сравнительный анализ алгоритмов}

В таблице~\ref{tab:comparison} приведены результаты сравнения подходов.

\begin{table}[H]
\centering
\caption{Сравнение решений, предложенных БЯМ}
\label{tab:comparison}
\begin{tabular}{|p{3.6cm}|p{3.6cm}|p{3.6cm}|p{3.6cm}|}
\hline
\textbf{Критерий} & \textbf{Gemini 3 Pro} & \textbf{Qwen 3 Coder (Local)} & \textbf{YandexGPT 5.1} \\ \hline
\textbf{Изначальная идея алгоритма} & Ошибочная (split) & Ошибочная (split) & \textbf{Верная} (поиск по raw-тексту) \\ \hline
\textbf{Работоспособность кода} & Высокая (запуск с 1 попытки) & Средняя (логические ошибки) & Низкая (ошибки исполнения) \\ \hline
\textbf{Реакция на замечания} & Быстрая самокоррекция & Требует уточнения формулировок & Игнорирование части замечаний \\ \hline
\end{tabular}
\end{table}

\subsection{Производительность и метрики генерации}

В ходе исследования зафиксированы существенные различия в поведенческих паттернах моделей, при этом общий объем сгенерированного текста (размер итоговых лог-файлов с диалогами) оказался практически идентичным для всех участников.

Основные наблюдения:
\begin{itemize}
    \item \textbf{временные затраты и режим ``размышления''} --- модель \textbf{Gemini 3 Pro} демонстрирует наибольшую задержку (latency), затрачивая около 60 секунд на этап скрытого планирования (Thinking process) и сопоставимое время на генерацию ответа. В противовес этому, \textbf{Qwen 3 Coder} и \textbf{YandexGPT} работают значительно быстрее, затрачивая на полный цикл ответа всего 20--30 секунд;
    \item \textbf{стиль коммуникации} --- локальная модель \textbf{Qwen 3} проявила себя как наиболее прагматичный инструмент, минимизирующий ``светскую беседу'' и пояснительную часть для выдачи преимущественно программного кода. Облачные модели, напротив, сопровождают решение развёрнутыми текстовыми комментариями и вступлениями;
    \item \textbf{итоговый объем данных} --- несмотря на кардинальные различия в стиле общения (``молчаливость'' Qwen против ``развернутости'' Gemini) и времени отклика, суммарный объем символьных данных в рамках сессии диалога оказался сопоставим.
\end{itemize}

\section{Выводы по исследовательскому разделу}

По результатам проведённого анализа сформулированы следующие выводы:
\begin{enumerate}
    \item \textbf{алгоритмическая интуиция YandexGPT} --- отечественная модель оказалась единственной, изначально верно спроектировавшей архитектуру поиска координат в тексте PDF (через подсчёт символов переноса строки), что позволило обойти ограничение библиотеки \texttt{PyPDF2}. Однако реализация этой идеи была осложнена грубыми синтаксическими ошибками, что характеризует модель как эффективный ``генератор идей'', но слабый инструмент кодирования без контроля со стороны человека;
    \item \textbf{стабильность Gemini 3 Pro} --- модель от Google продемонстрировала наилучший баланс. Несмотря на изначальный выбор более простого (и ошибочного для данной задачи) пути разбиения строк, она проявила высокую способность к исправлению и предоставила наиболее чистый и поддерживаемый код;
    \item \textbf{особенности локальной Qwen 3} --- локальная модель успешно справилась с задачей, однако потребовала от оператора выполнения роли ``ведущего''. Аналогично YandexGPT, модель с трудом интерпретировала запрос на разделение условий для предлогов ``от'' и ``до'', требуя максимально формализованного технического задания.
\end{enumerate}

\section{Анализ стабильности генерации и влияния режима рассуждений}

С целью исключения фактора случайности (<<выбросов>>) и оценки стабильности предлагаемых архитектурных решений проведена серия контрольных запусков.

\subsection{Методика исследования}
Каждой конфигурации модели четырежды (4 итерации) отправлялся идентичный первичный запрос. Анализу подвергался исключительно \textbf{первый ответ} модели.

Цель исследования заключалась в выяснении следующих аспектов:
\begin{itemize}
    \item влияет ли активация режима рассуждений (Reasoning/Thinking) на выбор стратегии решения (например, выбор между ошибочным методом \texttt{.split()} и верным поиском по индексам);
    \item насколько вариативен (стабилен) ответ модели от запуска к запуску.
\end{itemize}

Исследование проводилось для следующих конфигураций:
\begin{enumerate}
    \item \textbf{Gemini 3 Pro Preview} (Thinking Level: High);
    \item \textbf{Gemini 3 Pro Preview} (Thinking Level: Low);
    \item \textbf{YandexGPT 5.1} (режим <<Рассуждать>> включён);
    \item \textbf{YandexGPT 5.1} (режим <<Рассуждать>> выключен);
    \item \textbf{Qwen 3 Coder} (Instruct Mode).
\end{enumerate}

\subsection{Gemini 3 Pro Preview (Thinking Level: High)}

В рамках серии из 4-х контрольных запусков с максимальным уровнем рассуждений (High Reasoning) (см. приложение~\ref{app:gemini_high_multirun}) модель продемонстрировала высокую семантическую стабильность, но выявила тенденцию к выбору стандартных алгоритмических паттернов.

\textbf{Анализ результатов:}
\begin{enumerate}
    \item \textbf{интерпретация задачи} --- во всех четырёх случаях (100\%) модель верно интерпретировала условие <<проверка отсутствия неразрывного пробела>>. Генерировались регулярные выражения, нацеленные на поиск \textit{ошибки} (обычного пробела \texttt{[ ]+} или \texttt{[ $\backslash t$]+}), а не на поиск корректных NBSP, что подтверждает эффективность этапа рассуждения при фильтрации логических ловушек в техническом задании;
    
    \item \textbf{алгоритмическая стагнация} --- несмотря на активацию глубокого мышления, во всех запусках была выбрана стратегия предварительного разбиения текста на строки (\texttt{lines = text.split(``$\backslash n$'')}). Данный факт подтверждает гипотезу о том, что без явного указания на граничные случаи (разрыв фразы переносом строки) модель склоняется к наиболее распространённому решению, даже если оно не покрывает 100\% ситуаций. Процесс <<мышления>> в данном случае был направлен на корректное составление регулярных выражений, а не на архитектуру парсинга PDF;
    
    \item \textbf{лингвистическая адаптивность} --- в одном из запусков был продемонстрирован выход за рамки формального технического задания посредством самостоятельной реализации учёта морфологии русского языка. Вместо жёсткого поиска слов <<рисунок>>, <<таблица>> были предложены паттерны \texttt{r''рисун[а-я]*''} и \texttt{r''таблиц[а-я]*''}. Это позволяет находить ошибки даже в склоняемых формах слов, что недоступно другим моделям в рамках данного исследования, и является прямым следствием этапа <<рассуждения>>, где анализировались возможные вариации входных данных.
\end{enumerate}

\textbf{Вердикт по конфигурации} --- режим High Thinking обеспечивает глубокое понимание бизнес-логики задачи, но не гарантирует выбор наиболее сложного инженерного решения (обход разрывов строк) без наводящего вопроса.

\subsection{Gemini 3 Pro Preview (Thinking Level: Low)}

В серии запусков с пониженным уровнем рассуждений (Low Reasoning) (см. приложение~\ref{app:gemini_low_multirun}) моделью продемонстрирована алгоритмическая стагнация (повторение одного и того же неоптимального паттерна) и склонность к нарушению ограничений технического задания.

\textbf{Анализ результатов:}
\begin{enumerate}
    \item \textbf{нарушение стека технологий} --- в одном из запусков проявлена <<излишняя инициатива>>. Несмотря на прямое требование в запросе использовать библиотеку \texttt{PyPDF2}, был сгенерирован код с приоритетным импортом библиотеки \texttt{pypdf}. Подобное поведение (отклонение от технического задания в пользу <<лучших практик>> по мнению модели) отличает режим Low от High, где инструкции соблюдались строго;
    
    \item \textbf{отсутствие вариативности алгоритма} --- во всех четырёх запусках (100\%) использована стратегия предварительного разбиения текста (\texttt{lines = text.split(``$\backslash n$'')}). Снижение уровня <<мышления>> не привело к появлению альтернативных идей, модель стабильно выбирает наиболее очевидный, но уязвимый для разорванных фраз метод разбиения по символу перевода строки.
\end{enumerate}

\textbf{Вердикт по конфигурации} --- режим Low Thinking делает модель менее предсказуемой в вопросах выбора инструментов (библиотек) и стиля кода, при этом не обеспечивая преимущества в качестве алгоритмов.

\subsection{YandexGPT 5.1 Pro (режим <<Рассуждать>> включён)}

Серия из 4-х контрольных запусков (см. приложение~\ref{app:yagpt_think_multirun}) продемонстрировала высокую синтаксическую стабильность модели в данной выборке. Это позволяет квалифицировать критическую ошибку области видимости (Scope Error), полученную в ходе первичного исследования при тех же настройках, как случайный статистический выброс (stochastic outlier).

\textbf{Анализ результатов:}
\begin{enumerate}
    \item \textbf{воспроизводимость рабочего кода} --- во всех четырёх файлах данной серии моделью сгенерирован валидный, исполняемый код с использованием актуальных методов \texttt{PyPDF2}. Это подтверждает способность модели генерировать синтаксически верные решения, однако наличие грубой ошибки в первичном тесте указывает на ненулевую вероятность сбоя генерации даже при включённом режиме рассуждений;

    \item \textbf{логическая инверсия задачи (Business Logic Failure)} --- главной проблемой модели является стабильное непонимание сути валидации. Модель успешно находит искомые словосочетания, но \textbf{игнорирует проверку типа пробела}:
    \begin{itemize}
        \item в запуске №1 осуществляется поиск символа \texttt{$\backslash$u00A0} (NBSP) с выводом его наличия, то есть алгоритм ищет норму, а не ошибку;
        \item в запусках №2, 3, 4 используются паттерны с метасимволом \texttt{$\backslash$s} (любой пробел), вследствие чего программа выводит список всех вхождений без разбора, не отделяя ошибочные обычные пробелы от корректных неразрывных;
    \end{itemize}

    \item \textbf{стагнация алгоритма} --- несмотря на <<рассуждения>>, алгоритм парсинга PDF остался примитивным во всех итерациях (\texttt{text.split(``$\backslash$n'')}). Модель не предприняла попыток улучшить точность координат использованием поиска по полному тексту страницы, как это наблюдалось в первичном тесте.
\end{enumerate}

\textbf{Вердикт по конфигурации} --- режим рассуждений в данной выборке обеспечил 100\% синтаксически рабочего кода, но оказался бессилен перед логической сложностью задачи. Модель функционирует как ``поисковик фраз'', а не как валидатор вёрстки.

\subsection{YandexGPT 5.1 Pro (режим <<Рассуждать>> выключен)}

Серия контрольных запусков с отключённым режимом рассуждений (см. приложение~\ref{app:yagpt_no_think_multirun}) продемонстрировала парадоксальный результат --- качество логики кода оказалось выше, чем в режиме с активированными рассуждениями.

\textbf{Анализ результатов:}
\begin{enumerate}
    \item \textbf{синтаксическая стабильность} --- во всех четырёх запусках (100\%) сгенерирован валидный, рабочий код. Аналогично предыдущей серии, модель корректно работала с библиотекой \texttt{PyPDF2}, не допуская галлюцинаций с методами или импортами;

    \item \textbf{улучшение бизнес-логики (Logical Superiority)} --- в отличие от конфигурации с рассуждениями, осуществлявшей простой поиск совпадений по регулярным выражениям, <<быстрая>> модель в 3-х из 4-х случаев (75\%) реализовала корректную пост-валидацию найденных вхождений:
    \begin{itemize}
        \item в двух реализациях добавлена явная проверка \texttt{if '$\backslash$xa0' not in found\_text}, что полностью соответствует требованию задания (найти отсутствие неразрывного пробела);
        \item в одной реализации внедрена проверка на наличие обычного пробела (\texttt{if ' ' in found\_text});
        \item лишь в одной реализации модель ограничилась поиском по регулярным выражениям без дополнительной проверки (аналогично поведению в режиме с включёнными рассуждениями).
    \end{itemize}
    Данный факт указывает на то, что базовые веса модели YandexGPT 5.1 лучше настроены на паттерны  программирования, чем её надстройка рассуждений;

    \item \textbf{стабильность алгоритма} --- архитектурно решения остались идентичными, применялось разбиение текста через \texttt{splitlines()} или \texttt{split(``$\backslash$n'')}.
\end{enumerate}

\textbf{Вердикт по конфигурации} --- отключение режима <<Рассуждать>> сделало модель более эффективной для данной задачи. Исключение этапа <<философствования>> над текстом позволило реализовать стандартную логику проверки строки, что привело к рабочему решению в большинстве запусков.

\subsection{Qwen3 Coder}

В серии контрольных запусков (см. приложение~\ref{app:qwen_multirun}) локальная модель продемонстрировала наиболее глубокие знания спецификаций Unicode и синтаксиса Python, однако показала \textbf{критически низкую стабильность} в интерпретации условий задачи (всего 25\% верных попаданий в логику).

\textbf{Анализ результатов:}
\begin{enumerate}
    \item \textbf{энциклопедический подход к спецсимволам} --- Qwen3 является единственной моделью, которая не ограничилась стандартным неразрывным пробелом (\texttt{$\backslash$u00A0}). В регулярных выражениях был использован диапазон \texttt{[$\backslash$u00A0$\backslash$u2000-$\backslash$u200A$\backslash$u202F$\backslash$u205F$\backslash$u3000]}, что свидетельствует об обучении на глубоко технической документации;

    \item \textbf{системная ошибка бизнес-логики} --- в 3 из 4 запусков модель \textbf{инвертировала задачу}, осуществляя поиск вхождений, где неразрывный пробел \textit{присутствует} (условие \texttt{if '$\backslash$xa0' in match}), и помечая их как нарушения. Это указывает на то, что модель плохо различает семантику ``отсутствие X'' и ``наличие X'' в запросе;

    \item \textbf{феномен ``самокоррекции через итерации''} --- единственный успешный случай примечателен тем, что модель, ``сомневаясь'' в правильности, сгенерировала последовательно три варианта функции:
    \begin{itemize}
        \item \texttt{find\_special\_combinations} (инвертированная логика);
        \item \texttt{find\_special\_combinations\_detailed} (промежуточная попытка);
        \item \texttt{find\_special\_combinations\_correct} (итоговая верная логика с проверкой \texttt{not in}).
    \end{itemize}
    Это уникальное поведение, при котором модель провела ``внутренний рефакторинг'' непосредственно в процессе генерации ответа, в некотором роде эмулируя режим размышления;

    \item \textbf{культура кода} --- код Qwen3 (аналогично YandexGPT Thinking Off) отличается высоким качеством оформления, включая активное использование аннотации типов (\texttt{List}, \texttt{Tuple}, \texttt{Optional}) и обработку исключений.
\end{enumerate}

\textbf{Вердикт по конфигурации} --- Qwen3 Coder проявила себя как ``технический эрудит'' с серьёзными проблемами в понимании отрицаний в техническом задании. Модель склонна к инверсии условий, и только генерация нескольких вариантов решения позволила ей определить верный путь.

\subsection{Сводные результаты серии запусков}

В таблице~\ref{tab:stability_test} представлены обобщённые метрики успешности первичной генерации кода. Под ``успешным'' понимается решение, которое синтаксически корректно и реализует верную бизнес-логику (поиск отсутствия NBSP).

\begin{table}[H]
\centering
\caption{Статистика успешности первичных решений (4 запуска)}
\label{tab:stability_test}
\begin{tabular}{|l|c|p{4cm}|l|}
\hline
\textbf{Модель / Режим} & \textbf{Успех} & \textbf{Типичная ошибка} & \textbf{Алгоритм} \\ \hline
Gemini 3 (High) & 4/4 & Неоптимальная архитектура & \texttt{split('$\backslash$n')} \\ \hline
Gemini 3 (Low) & 3/4 & Галлюцинации (libs), стиль & \texttt{split('$\backslash$n')} \\ \hline
YaGPT 5.1 (On) & 0/4 & Инверсия логики (ищет норму) & \texttt{split('$\backslash$n')} \\ \hline
YaGPT 5.1 (Off) & 3/4 & Инверсия логики (редко) & \texttt{splitlines()} \\ \hline
Qwen 3 Coder & 1/4 & Инверсия логики (часто) & \texttt{split('$\backslash$n')} \\ \hline
\end{tabular}
\end{table}

\section*{Выводы}

Исследованием выявлено фундаментальное различие в том, как ``рассуждающие'' и стандартные модели обрабатывают отрицания в технических заданиях (проверка \textit{отсутствия} элемента).

\begin{enumerate}
    \item \textbf{парадокс YandexGPT} --- для отечественной модели активация режима рассуждений (CoT) оказала негативное влияние на результат. ``Думающая'' версия стабильно игнорировала суть валидации, работая как поисковик, тогда как базовая версия (Thinking Off) успешно реализовывала запрашиваемую логику в 75\% случаев, что свидетельствует о переобучении модуля рассуждений на задачи поиска информации, а не анализа кода;
    
    \item \textbf{семантическое превосходство Gemini (High)} --- единственной конфигурацией, обеспечившей 100\% попадание в бизнес-логику задачи, стала Gemini 3 с максимальным уровнем рассуждений. Модель безошибочно распознала требование искать ``ошибки'', а не ``правильные варианты'', несмотря на использование стандартного алгоритма парсинга;
    
    \item \textbf{проблема отрицаний у Code-моделей} --- специализированные модели (Qwen 3 Coder) и YandexGPT (Thinking On) склонны к ``инверсии задачи''. Фраза ``проверка отсутствия'' часто интерпретируется ими как ``проверка наличия'', что представляет критический риск при автоматической генерации тестов и валидаторов;
    
    \item \textbf{итоговая рекомендация} --- для задач, требующих сложной логической валидации и работы с отрицаниями, необходимо использовать SOTA-модели с глубоким рассуждением (Gemini High). Для генерации синтаксически сложного, типизированного кода (boilerplate) эффективнее использовать локальные или быстрые модели (Qwen, YaGPT Off), но с обязательным ручным ревью бизнес-логики.
\end{enumerate}

\clearpage